\section{案例研究：现实世界的 Offline RL}

\subsection{复杂 reward 工程}

在人机交互系统中，reward often包含多个目标（accuracy, fairness, latency）。讲者描述 LinkedIn 项目：在 notification 系统中同时优化 CTR 和 retention，因此要在 dataset 中标记 reward components 以便 later decomposed training。

\begin{importantbox}{多目标 reward 的 practical trick}
在 dataset 中预先计算 reward vector（CTR, retention, quality），并在 offline training 时用 linear combination + CVaR regularization 控制 worst-case outcomes。
\end{importantbox}

\subsection{可视化演示}

\begin{figure}[H]
\centering
\includegraphics[page=6,width=0.9\textwidth]{07_cs224r_offline_rl_2025.pdf}
\caption{Slides 表示的 reality check：多目标 reward 与 deployment 指标对齐\protect\footnotemark}
\end{figure}
\footnotetext{Slides 第6页，强调 reward vector 与 deployment metrics 的一一映射。}

\subsection{本章小结}

现实案例中 reward 是多维的，务必在数据层面明确每个维度的意义，否则 offline policy 会 misalign。

